\section{End-Term Comprehensive Examination --- Set B}
\label{sec:endterm_set_b}

\LPUPaperHeader{End-Term Comprehensive Examination}{B}{180 Minutes (3 Hours)}{50 Marks}{CO1--CO6 (Comprehensive Syllabus Units 1 to 6)}

\subsection*{Section A: Objective Multiple Choice Questions [10 Marks]}
\begin{center}
\sffamily\textbf{\color{AlertRed}Attempt all ten questions. Each question carries 1 mark.}
\end{center}

\mcqitem{1}{Which framework categorizes decision contexts into Clear, Complicated, Complex, and Chaotic?}{Double Diamond}{Cynefin Framework}{MCDA Matrix}{Ishikawa Model}{1}

\mcqitem{2}{What is the primary formula structure for a human-centered Point-of-View (POV) statement?}{If-Then-Else}{[User] needs [Need] because [Surprising Insight]}{Input $\to$ Process $\to$ Output}{Problem = Symptom + Cost}{1}

\mcqitem{3}{Which diagnostic tool explores causal chains by repeatedly asking 'Why'?}{Fishbone 6M}{The 5-Whys Method}{Pugh Matrix}{Affinity Diagram}{1}

\mcqitem{4}{What does the 'E' in SCAMPER stand for?}{Execute}{Evaluate}{Eliminate}{Expand}{1}

\mcqitem{5}{In Stuart Pugh's Controlled Convergence Matrix, how is the baseline datum scored across all criteria?}{Always +1}{Always 0}{Always -1}{Always 10}{1}

\mcqitem{6}{What does a System Usability Scale (SUS) score of 50.0 represent?}{Top 5\% usability}{Exact neutral response across all items}{Perfect system usability}{Zero usability}{1}

\mcqitem{7}{What is the boundary between customer-visible actions and backstage operations in a Service Blueprint?}{Line of Interaction}{Line of Visibility}{Line of Internal Coordination}{Line of Separation}{1}

\mcqitem{8}{What are the four components of Barbara Minto's SCQA framework?}{Speed, Cost, Quality, Action}{Situation, Complication, Question, Answer}{Scope, Criteria, Quantification, Assessment}{Strategy, Concept, Query, Architecture}{1}

\mcqitem{9}{Which stage of Graham Gibbs' Reflective Cycle asks 'What else could you have done?'}{Description}{Analysis}{Conclusion}{Feelings}{1}

\mcqitem{10}{What is the primary principle of blameless engineering postmortems?}{Assigning individual fault quickly}{Treating human error as a symptom of deeper system design flaws}{Issuing monetary fines}{Ignoring system crashes}{1}

\newpage
\subsection*{Section B: Applied Analytical Engineering Scenarios [20 Marks]}
\begin{center}
\sffamily\textbf{\color{AlertRed}Attempt all four questions. Each question carries 5 marks.}
\end{center}

\questionitem{11}{
    A university data center is selecting an enterprise storage architecture from three alternatives: SAN, NAS, and Distributed Object Storage. The criteria and weights are: Read/Write Throughput ($w_1 = 0.30$), Scalability ($w_2 = 0.30$), Disaster Recovery Redundancy ($w_3 = 0.25$), and Total Cost of Ownership ($w_4 = 0.15$). The performance ratings (1--10) are:
    \begin{itemize}
        \item \textbf{Option 1 (SAN):} Throughput = 10, Scalability = 6, Redundancy = 8, Cost = 5
        \item \textbf{Option 2 (NAS):} Throughput = 7, Scalability = 8, Redundancy = 7, Cost = 8
        \item \textbf{Option 3 (Distributed Object):} Throughput = 8, Scalability = 10, Redundancy = 9, Cost = 7
    \end{itemize}
    Compute the MCDA composite scores and identify the winning architecture.
}{5}{CO3}{L3 Apply}

\questionitem{12}{
    Explain the difference between Inductive and Abductive reasoning in the context of user research and product feature discovery.
}{5}{CO3}{L2 Understand}

\questionitem{13}{
    Design an MVP validation experiment for an on-demand drone organ transport system between regional hospitals. Choose between a Smoke Test, Wizard of Oz, or Concierge MVP and justify your selection.
}{5}{CO5}{L4 Analyze}

\questionitem{14}{
    Construct a 2$\times$2 Impact vs. Effort Prioritization Matrix classifying four potential software features for an automated student examination proctoring tool.
}{5}{CO4}{L3 Apply}

\newpage
\subsection*{Section C: Comprehensive Capstone Design Challenges [20 Marks]}
\begin{center}
\sffamily\textbf{\color{AlertRed}Attempt both comprehensive design challenges. Each question carries 10 marks.}
\end{center}

\questionitem{15}{
    \textbf{Design Challenge 1: Smart Campus Energy Conservation}\\
    A university campus with 35,000 students experiences massive electricity waste: lecture halls have air conditioning and high-intensity lights running continuously in empty buildings, wasting \$850,000 annually. The administration considered installing manual timer switches, but faculty and staff complained they were disruptive.\\
    \textbf{Deliverables Required:}\\
    (a) Structure a root cause analysis using both 5-Whys and an Ishikawa diagram (3 Marks).\\
    (b) Formulate an empathetic POV statement and derive three actionable HMW questions (3 Marks).\\
    (c) Apply SCAMPER to generate four creative smart building interventions (4 Marks).
}{10}{CO2, CO4}{L5 Evaluate}

\questionitem{16}{
    \textbf{Design Challenge 2: Service Blueprinting \& Blameless Postmortem}\\
    An automated campus drone meal delivery service launched during mid-terms crashed spectacularly: 45 drones collided with trees, 200 orders were delivered cold or dropped onto car roofs, and students flooded social media with mocking videos.\\
    \textbf{Deliverables Required:}\\
    (a) Design a complete Multi-Tier Service Blueprint identifying lines of interaction, visibility, and key failure nodes [!] (5 Marks).\\
    (b) Conduct a comprehensive Blameless Postmortem detailing timeline reconstruction, systemic contributing factors, and a concrete remediation action plan (5 Marks).
}{10}{CO5, CO6}{L5 Evaluate}

\newpage
\subsection*{Solutions and Marking Rubrics --- End-Term Set B}

\begin{solutionbox}[Solutions to Section A: Objective MCQs]
\textbf{1.} (B) Cynefin Framework \quad
\textbf{2.} (B) [User] needs [Need] because [Surprising Insight] \quad
\textbf{3.} (B) The 5-Whys Method \quad
\textbf{4.} (C) Eliminate \quad
\textbf{5.} (B) Always 0 \quad
\textbf{6.} (B) Exact neutral response across all items \quad
\textbf{7.} (B) Line of Visibility \quad
\textbf{8.} (B) Situation, Complication, Question, Answer \quad
\textbf{9.} (C) Conclusion \quad
\textbf{10.} (B) Treating human error as a symptom of deeper system design flaws
\begin{markingrubric}
1 Mark per correct MCQ = \textbf{10 Marks Total}.
\end{markingrubric}
\end{solutionbox}

\begin{solutionbox}[Solution to Question 11: Data Center Storage MCDA]
Calculate $S_i = \sum w_j \cdot r_{ij}$:\\
$\bullet$ \textbf{Option 1 (SAN):}\\
$S_1 = (0.30 \times 10) + (0.30 \times 6) + (0.25 \times 8) + (0.15 \times 5) = 3.00 + 1.80 + 2.00 + 0.75 = \mathbf{7.55}$\\
$\bullet$ \textbf{Option 2 (NAS):}\\
$S_2 = (0.30 \times 7) + (0.30 \times 8) + (0.25 \times 7) + (0.15 \times 8) = 2.10 + 2.40 + 1.75 + 1.20 = \mathbf{7.45}$\\
$\bullet$ \textbf{Option 3 (Distributed Object):}\\
$S_3 = (0.30 \times 8) + (0.30 \times 10) + (0.25 \times 9) + (0.15 \times 7) = 2.40 + 3.00 + 2.25 + 1.05 = \mathbf{8.70}$\\
\textbf{Recommendation:}\\
Distributed Object Storage emerges as the definitive winner ($S_3 = 8.70$), outperforming SAN and NAS due to near-infinite scalability, robust built-in erasure coding replication, and low commodity storage costs.
\begin{markingrubric}
$\bullet$ Accurate calculation across all three candidates: \textbf{3 Marks}\\
$\bullet$ Clear analytical justification and final selection: \textbf{2 Marks}
\end{markingrubric}
\end{solutionbox}

\begin{solutionbox}[Solution to Question 12: Inductive vs. Abductive in Design]
$\bullet$ \textbf{Inductive Reasoning (Pattern Extraction):}\\
Observes repeated behavioral occurrences to infer general statistical truths. Example: User telemetry shows that 74\% of mobile app users drop off when asked to type their address manually $\implies$ Inductive conclusion: Manual forms induce friction.\\
$\bullet$ \textbf{Abductive Reasoning (Inference to Best Explanation):}\\
Observes an anomalous, surprising behavior and makes an empathetic creative leap to hypothesize the latent root cause. Example: An older user pauses for 25 seconds on a 'Submit' button despite zero errors displayed. An inductive approach sees only delay. An abductive leap realizes: ``The user fears that clicking will irrevocably charge their credit card without a confirmation step.'' This hypothesis directly generates a confirmation modal prototype.
\begin{markingrubric}
$\bullet$ Clear epistemological distinction: \textbf{2 Marks}\\
$\bullet$ Specific, highly insightful design/user research illustrations: \textbf{3 Marks}
\end{markingrubric}
\end{solutionbox}

\begin{solutionbox}[Solution to Question 13: Drone Organ Transport MVP Selection]
\textbf{Optimal Strategy: High-Fidelity Prototype Spike with Human-in-the-Loop Shadowing.}\\
$\bullet$ \textit{Why Smoke Test / Fake Door is Inappropriate:} High ethical stakes and patient life dependency make deceptive mock landing pages unethical and legally prohibited in organ logistics.\\
$\bullet$ \textit{Validation Architecture:} Deploy a single certified drone on a pre-cleared 4-kilometer corridor between two partner teaching hospitals to transport non-biological dummy payload bags with thermal sensors and GPS loggers.\\
$\bullet$ \textit{Critical Hypotheses Tested:} Vibration levels, temperature maintenance ($\pm 0.5^\circ\text{C}$), landing precision in wind gusts, and handover latency between flight technician and surgical nurse.
\begin{markingrubric}
$\bullet$ Rigorous critique of why simple smoke tests fail in high-risk contexts: \textbf{2 Marks}\\
$\bullet$ Detailed, actionable, ethical validation experiment protocol: \textbf{3 Marks}
\end{markingrubric}
\end{solutionbox}

\begin{solutionbox}[Solution to Question 14: Proctoring Tool 2$\times$2 Prioritization]
\textbf{Classification of Proctoring Features:}\\
1. \textbf{Quick Wins (High Impact, Low Effort):} Instant multi-tab browser blur detection alert.\\
2. \textbf{Strategic Bets (High Impact, High Effort):} Low-bandwidth edge AI gaze-tracking algorithm that runs on client CPU without sending raw video feeds to the cloud.\\
3. \textbf{Fill-ins (Low Impact, Low Effort):} Customizable dark mode for the exam interface.\\
4. \textbf{Money Pits (Low Impact, High Effort):} Continuous biometric 3D facial mesh reconstruction via webcams (prohibitive client CPU load, high false rejections, severe student privacy backlash).
\begin{markingrubric}
$\bullet$ Correct placement across all 4 quadrants: \textbf{2 Marks}\\
$\bullet$ Clear technical and pedagogical justification: \textbf{3 Marks}
\end{markingrubric}
\end{solutionbox}

\begin{solutionbox}[Solution to Question 15: Campus Energy Conservation Capstone]
\textbf{(a) Root Cause Analysis (3 Marks):}\\
$\bullet$ \textit{5-Whys:} Why are empty rooms cooled? Because AC runs on an unmonitored master building schedule $\to$ Why master schedule? Because individual room occupancy is not tracked $\to$ Why not tracked? Because room sensors were deemed too costly in 2010 $\to$ Why do faculty not switch them off? Because switches are locked in electrical closets to prevent tampering $\to$ \textit{Root Cause:} Centralized rigid scheduling decoupled from real-time class registration schedules.\\
$\bullet$ \textit{Ishikawa Breakdown:} People (faculty lack awareness), Policy (facilities department locks master panels), Technology (no dynamic room booking API integration).\\
\textbf{(b) POV Statement \& HMW Questions (3 Marks):}\\
$\bullet$ \textit{POV:} ``University facilities managers [User] need automated, occupant-aware ambient climate control [Need] because rigid centralized schedules and lack of room-level feedback waste massive energy and budget while frustrating occupants [Insight].''\\
$\bullet$ \textit{HMW 1:} How might we link classroom AC power directly to the live registrar course schedule database?\\
$\bullet$ \textit{HMW 2:} How might we empower students to 'check into' a room to activate cooling for study sessions?\\
$\bullet$ \textit{HMW 3:} How might we gamify energy savings between campus departments through transparent floor dashboards?\\
\textbf{(c) SCAMPER Solutions (4 Marks):}\\
$\bullet$ \textbf{Substitute:} Substitute mechanical timers with calendar-integrated smart IoT thermostat relays.\\
$\bullet$ \textbf{Combine:} Combine Wi-Fi access point device count telemetry with HVAC airflow throttling.\\
$\bullet$ \textbf{Modify:} Modify cooling set-points: drift temperature dynamically by $1.5^\circ\text{C}$ during 15-minute passing periods.\\
$\bullet$ \textbf{Eliminate:} Eliminate manual master electrical closet overrides entirely.
\begin{markingrubric}
$\bullet$ Part (a) 5-Whys and Ishikawa analysis: \textbf{3 Marks}\\
$\bullet$ Part (b) POV statement and 3 HMW questions: \textbf{3 Marks}\\
$\bullet$ Part (c) 4 SCAMPER solutions: \textbf{4 Marks}
\end{markingrubric}
\end{solutionbox}

\begin{solutionbox}[Solution to Question 16: Drone Crash Blueprint \& Blameless Postmortem]
\textbf{(a) Service Blueprint Design (5 Marks):}\\
$\bullet$ \textit{Physical Evidence:} Mobile meal tracking app; insulated drone tether box; thermal bag; notification ping.\\
$\bullet$ \textit{Customer Actions:} Orders burger $\to$ walks to outdoor campus drop circle $\to$ confirms arrival pin $\to$ unlatches box.\\
$\bullet$ \textit{Line of Interaction:}\\
$\bullet$ \textit{Frontstage Contact:} Drone descends to 3-meter hover, lowers tethered box into painted yellow circle.\\
$\bullet$ \textit{Line of Visibility:}\\
$\bullet$ \textit{Backstage Contact:} Restaurant kitchen packs hot meal into aerodynamically balanced container; ground crew snaps container onto drone hook.\\
$\bullet$ \textit{Line of Internal Interaction:}\\
$\bullet$ \textit{Support Processes:} Autonomous flight collision avoidance mesh; LIDAR wind-gust compensator; campus airspace traffic scheduler.\\
$\bullet$ \textit{Failure Nodes [!]:} [! 1] High wind gusts destabilizing tether winch; [! 2] Tree canopy obscuring visual landing fiducial marker.\\
\textbf{(b) Blameless Postmortem Protocol (5 Marks):}\\
1. \textit{Incident Timeline:} 12:15 PM: 200 orders surge simultaneously; 12:22 PM: GPS signal reflection off campus glass tower causes 45 drones to drift into trees; 12:35 PM: Emergency manual override server crashes under load.\\
2. \textit{Systemic Contributing Factors:} Gaps in multi-path GPS filtering near glass buildings; absence of automated winch cut-off on tether snag; failure to stress-test concurrent peak meal volume in simulation.\\
3. \textit{Corrective Remediation Plan:} P0 Ticket \#401: Install RTK ground base stations for centimeter-accurate positioning (Due in 7 days); P0 Ticket \#402: Implement failsafe auto-abort climb when tether tension exceeds 3 kg (Due in 5 days); P1 Ticket \#403: Cap concurrent delivery slots at 25 drones per 15-minute window.
\begin{markingrubric}
$\bullet$ Part (a) Multi-tier blueprint with failure nodes: \textbf{5 Marks}\\
$\bullet$ Part (b) Blameless postmortem with timeline, root causes, and P0/P1 remediation tickets: \textbf{5 Marks}
\end{markingrubric}
\end{solutionbox}
